\documentclass[../main.tex]{subfiles}
\begin{document}
%%%%%%%%%%%%%%%%%%%%%%%%%%%%%%%%%%%%%%%%%%%%%%%%%%%%%%%%%%%%%%%%%%%%%%%%%%%%%%%%%%%%%%%%%%%%%%%%%%%
\chapter{Exporting Image and Graphics Files from OmniGraffle}

One of the main reasons you use a design app like OmniGraffle is so you can export your work as
image or graphics files for use with other projects. OmniGraffle exports files to the standard
image types, such as PNG, JPEG, and TIFF, as well as to PDF and EPS. And if you have
OmniGraffle Pro, you can export files to Microsoft Visio and as SVG (Scalable Vector
Graphics).

This chapter introduces you to OmniGraffle's new \textbf{Export} panel, from which you can
select the file type to export, specify multiple image resolutions to output (such as @2x
through @5x for web and app development), and more.

\section{Exportable File Types}

The following table serves as a quick reference for the image and file types that you can export
from OmniGraffle.

\begingroup
\small
\setlength{\LTleft}{0.075\textwidth}
\setlength{\LTright}{0.075\textwidth}
\renewcommand{\arraystretch}{1.25}
\begin{xltabular}{0.85\textwidth}{|>{\raggedright\arraybackslash}X|c|c|}
  \caption{List of exportable file types}\\
  \hline
  \rowcolor{ogheaderblue}
  \textcolor{white}{\textbf{File Type}} & \textcolor{white}{\textbf{Standard}} &
  \textcolor{white}{\textbf{Pro}} \\
  \hline
  \endhead
  \hline
  \endfoot
  PNG (\textit{.png}) & \ogyes & \ogyes \\
  \hline
  JPEG (\textit{.jpg}, \textit{.jpeg}) & \ogyes & \ogyes \\
  \hline
  GIF (\textit{.gif}) & \ogyes & \ogyes \\
  \hline
  BMP (\textit{.bmp}) & \ogyes & \ogyes \\
  \hline
  TIFF (\textit{.tiff}) & \ogyes & \ogyes \\
  \hline
  Adobe Photoshop (\textit{.psd}) & \ogyes & \ogyes \\
  \hline
  OmniOutliner (\textit{.oo3}) & \ogyes & \ogyes \\
  \hline
  OmniGraffle Document (\textit{.graffle}) & \ogyes & \ogyes \\
  \hline
  OmniGraffle Stencil (\textit{.gstencil}) & \ogyes & \ogyes \\
  \hline
  OmniGraffle Template (\textit{.gtemplate}) & \ogyes & \ogyes \\
  \hline
  PDF (\textit{.pdf}) & \ogyes & \ogyes \\
  \hline
  EPS (\textit{.eps}) & \ogyes & \ogyes \\
  \hline
  SVG (\textit{.svg}) & \ogno & \ogyes \\
  \hline
  Microsoft Visio (\textit{.vdx}) & \ogno & \ogyes \\
  \hline
\end{xltabular}
\endgroup

\section{Using the Export Panel}

You can create files of several different types from your OmniGraffle projects. Choose
\textbf{File} $\blacktriangleright$ \textbf{Export} (\textbf{Option-Command-E}) to open the
Export panel:

\ogfigure{0.800}{res/og7_export01.png}

The basic steps of exporting from OmniGraffle are as follows:

\begin{enumerate}
  \item Choose an export file type from the row along the top of the export panel.
  \item Define what gets exported on the top-left portion of the panel.

    \begin{ogtip}
      If you have set a \textbf{Canvas Fill} or \textbf{Background Image} (see Working With the
      Inspectors) that you would like included as part of the export area, turn off
      \textbf{Transparent background} in the Export panel.
    \end{ogtip}
  \item Add additional \textbf{Export Sizes} to the list, if necessary.
  \item Examine the \textbf{Export Preview} to make sure that what gets exported meets your
    expectations.
  \item Click \textbf{Export}.
\end{enumerate}

After clicking \textbf{Export}, a Save sheet appears over the export panel, from which you can
choose where the exported file (or files, if you chose extra Export Sizes) is saved. By
default, OmniGraffle uses the name of your OmniGraffle project as the filename---plus any
\textbf{Name Suffix} you've added in the \textbf{Export Sizes} area---and then tacks on the
appropriate file extension (for example, \textit{MyAppIcon.png} and
\textit{MyAppIcon@2x.png}).

Of course, there's much more to the export panel than meets the eye, so let's take a closer look
at all of the things you can set and choose from in the panel.

\section{Defining Your Export in the Export Panel}

The export panel is divided into four areas:

\begin{itemize}
  \item the \textbf{Export File Type} row along the top;
  \item the \textbf{Export Definition} area at top-left;
  \item the \textbf{Export Sizes} area at bottom-left; and,
  \item the \textbf{Export Preview} area on the right side of the panel.
\end{itemize}

Let's dive a little deeper to see what's all in there\ldots{}

\subsection{Choosing an Export File Type}

The row at the top of the export panel includes clickable icons that you use to select the type
of file you want to export. OmniGraffle Pro users also have the ability to export files in
SVG, Microsoft Visio, and HTML Image Map formats.

\ogfigure{0.800}{res/og7_export02.png}

The Export panel remembers the file type of your previous export, and auto-selects that file
format in the row along the top. This makes it easy for you to repeat exports from the same or
different OmniGraffle projects. If you need to export to a different file type, click to
choose the type you need.

\subsection{Export Definition}

The top-left area of the export panel (everything above the \textbf{Export Sizes} area) help you
define what gets exported. This area changes based on the export file type you have selected
in the top row; for example, this is how the definition area appears with PNG selected:

\ogfigure{0.800}{res/og7_export03.png}

The \textbf{Format} section lists the type of file to be exported, along with its file
extension. If the image type you have selected supports transparency, you'll also see a
\textbf{Transparent background} option, which you can turn on or off, depending on your needs.

\begin{ognote}
  The JPEG and Windows Bitmap (BMP) file formats do not offer transparency. Nor do the
  OmniOutliner and OmniGraffle file types.
\end{ognote}

The \textbf{Export area} pop-up menu offers the following options for you to choose how much of
document to export:

\begin{itemize}
  \item \textbf{Entire Document} --- Export everything, creating one image file for each canvas.
  \item \textbf{Current Canvas} --- Export the entire canvas.
  \item \textbf{All Objects (Current Canvas)} --- Export all of the objects on the current
    canvas. When you choose \textbf{All Objects (Current Canvas)}, every object on that canvas
    is included as part of the export set. If an object appears outside the bounds of the canvas
    area, it too will be included as part of the export.
  \item \textbf{Selection (Current Canvas)} --- Only export the objects that are selected right
    now.
\end{itemize}

By default, non-printing layers are excluded from being exported; however, for most image file
types, such as PNG and JPEG, you also have the option to \textbf{Include non-printing layers}
as part of the export. For example, you can turn off printability for layers in the Sidebar by
clicking the \ogicon{12.5pt}{res/og7_icon_layer_print.png}, flipping it to
\ogicon{12.5pt}{res/og7_icon_layer_dontprint.png} to disable printing for that layer. If you
want to include non-printing layers as part of the export, go ahead and turn on this option.

If your project uses artboards and you want to export images based on the artboard shapes, make
sure that \textbf{Export from artboards only} is checked. This tells OmniGraffle to export
images based on the artboards in your project; without this option checked, what gets exported
depends on what you have selected in the Export Area pop-up menu.

Available for PNG, JPEG, GIF, BMP, and TIFF image types, you'll also see \textbf{Resolution}
section. By default, all image types are exported at \textbf{72 dots per inch}, but you can
change this to whatever you'd like. If you prefer, you can change \textbf{dots per inch} to
\textbf{dots per cm} in the pop-up menu.

JPEG exports also offer a \textbf{Quality} slider, which you can use to define how lossy you
want those images to be.

\ogfigure{0.800}{res/og7_export04.png}

While EPS export lets you choose whether the exported file has a preview image:

\ogfigure{0.800}{res/og7_export05.png}

\subsection{Export Sizes}

For most image formats, you can choose to export in varying sizes. The default \textbf{Size} is
100\%, based on the actual size of the things you are exporting, and can range up to 500\% or
as low as 1\%, depending on your design needs.

\ogfigure{0.800}{res/og7_export06.png}

To add an additional export size, click \ogicon{12.5pt}{res/og7_button_export_addsize.png} next
to the first export size option. OmniGraffle increments the \textbf{Size} by another 100\%
each time you click that button; however, you can specify any size variant you'd like, as well
as the \textbf{Name Suffix} that gets added to the images you export.

If you are only exporting the natural size (\textbf{100\%}), the filename given to the exported
file matches the filename of your OmniGraffle project. That said, you can change that filename
to whatever you'd like after clicking the \textbf{Export} button.

If you are exporting multiple sizes, however, the filename is based on the canvas name you've
specified in the Sidebar. Additionally, after clicking the \textbf{Export} button, you'll
notice that the \textbf{Export As} field uses the filename for your OmniGraffle project. The
filename, which you can change, is actually a folder that gets created in a location you
specify for containing the exported images. When you click \textbf{Save}, OmniGraffle places
all of the exported image files into that folder.

\ogfigure{0.800}{res/og7_export07.png}

\ogfigure{0.800}{res/og7_export08.png}

\subsection{Export Preview}

The Export Preview takes up the entire right side of the panel. If you are exporting multiple
sizes, or from multiple canvases in your project, the preview area includes right and left
arrow buttons so you can page through and preview the images in your export set.

\ogfigure{0.800}{res/og7_export09.png}

At the top of the preview, you'll see that OmniGraffle displays the intended filename for that
image, its size in width and height, as well as the final exported file size.

\section{Exportable Omni File Types}

In addition to exporting to various image types, you can also export your OmniGraffle projects
as OmniOutliner and OmniGraffle files.

If you use OmniOutliner, one reason why you might want to export to that format is if you are
working on a diagram, such as an organization chart or process workflow.

\ogfigure{0.800}{res/og7_export10.png}

Exporting to OmniOutliner works best for diagrams that adhere to a tree-like structure.
OmniGraffle uses the connection lines between shapes to create the outline hierarchy.

When exporting as an OmniGraffle file, choose the \textbf{Format} pop-up menu to export your
project as another OmniGraffle file (\textit{.graffle}), as an OmniGraffle Stencil
(\textit{.gstencil}), or as an OmniGraffle Template (\textit{.gtemplate}).

\ogfigure{0.800}{res/og7_export11.png}

Use the second pop-up menu under \textbf{Format} to choose the OmniGraffle file type for export
from among the following:

\begin{itemize}
  \item \textbf{Single File (Zipped)}---The default file type for all files created with
    OmniGraffle 7.16 and later, best suited for cloud storage.
  \item \textbf{Package}---An alternative file type with some niche uses for improving document
    performance. Not recommended for cloud storage.
  \item \textbf{Legacy (Uncompressed Single File)}---A file type compatible with previous
    versions of OmniGraffle, replaced by the faster zipped single file type. Suited for cloud
    storage when syncing with older versions.
  \item \textbf{Legacy (Compressed Single File)}---A file type compatible with previous versions
    of OmniGraffle, suited for local archival storage. Not recommended for any sort of cloud
    storage, version control, or scripting.
  \item \textbf{Legacy (Uncompressed Package)}---A file type compatible with previous versions
    of OmniGraffle, suited for better document performance in some circumstances. Not
    recommended for cloud storage.
\end{itemize}

Use the checkboxes to control aspects of the exported file's size and editability. You can opt
to select from the following:

\begin{itemize}
  \item \textbf{Copy linked images into document} is useful for when your OmniGraffle project
    has image files that you've added to a shape or the canvas and you'd like those to go along
    with the file.
  \item \textbf{Make documents read-only} is particularly helpful for stencils that you are
    sharing with a design team and you'd prefer that they don't edit the objects.
\end{itemize}

\section{Exporting an HTML Image Map (\texorpdfstring{\textbf{PRO}}{PRO})}

\begin{ogpro}
  HTML Image Maps have been around for a long time, and continue to be a useful interface element
  on the web. HTML image maps typically appear as single, cohesive image on a website that, when
  moused over, present users with the option to click an area and then go to another part of the
  website.

  \begin{ognote}
    Back in the day, there were no great tools for creating HTML image maps; you had to work out
    the coordinates and define the click areas on your own. Thankfully, OmniGraffle takes that
    pain away so you can concentrate on your design work.
  \end{ognote}

  HTML image maps are unlike any other image export you'll do in OmniGraffle. The reason is
  because they leverage the power of Actions, which you use to define URLs for the clickable
  regions of your design. Also, OmniGraffle exports the HTML to go along with the images so
  you can upload the assets to your website and instantly put them to use.

  For example, let's say that you needed to create a navbar for your website. You can design
  this in OmniGraffle, as shown here:

  \ogfigure{0.800}{res/og7_export12.png}

  The navbar has a ``Home'' button on the left side, along with four text areas for links that
  take you deeper into the website. You would use the \textbf{Action} inspector to assign
  links to each clickable item in the navbar; for example, here's how the ``Home'' button was
  configured:

  \ogfigure{0.800}{res/og7_export13.png}

  When you're ready to export the files, choose \textbf{File} $\blacktriangleright$
  \textbf{Export} (\textbf{Option-Command-E}) and then select \textbf{HTML Image Map} in the
  filetype row at the top of the export panel:

  \ogfigure{0.800}{res/og7_export14.png}

  When you export the files, OmniGraffle generates an HTML file and an \textbf{HTML Image Type},
  based on your choice on the left side of the panel. URL actions in the original OmniGraffle
  document are coded into the image map as links so that the image can be clicked to follow
  them. (Note that actions on lines aren't included, but actions on line labels are.)

  The source files are saved to a folder, which you can then test and integrate with your
  website design.
\end{ogpro}
%%%%%%%%%%%%%%%%%%%%%%%%%%%%%%%%%%%%%%%%%%%%%%%%%%%%%%%%%%%%%%%%%%%%%%%%%%%%%%%%%%%%%%%%%%%%%%%%%%%
\end{document}
